\section{A common integral gluing complex}\label{sec:common-chain}

\subsection{The common period marking}

Let $U=\mathbb P^1\setminus\{0,1,\infty\}$, with its complex orientation.
For composable paths, $\alpha\beta$ means first traverse $\beta$, then
traverse $\alpha$. Path multiplication is therefore right to left.
Choose a base point and marking paths for counterclockwise meridians
such that
\begin{equation}\label{cc:eq:path-order}
\gamma_0\gamma_1\gamma_\infty=1,
\qquad \gamma_\infty=(\gamma_0\gamma_1)^{-1}
=\gamma_1^{-1}\gamma_0^{-1}.
\end{equation}
The first word traverses $\gamma_\infty$, then $\gamma_1$, then $\gamma_0$.
All fundamental groups, transport maps, and presentations use this convention.

The fiber is the real torus $F=\mathbb R^4/\Lambda$, where
$\Lambda=\mathbb Ze_1\oplus\cdots\oplus\mathbb Ze_4$ has its displayed orientation.
The two complex period markings are
\[
\begin{aligned}
e_1&=(1,0),&e_2&=(\rho,0),&e_3&=(2+\rho,1)/3,&e_4&=(0,\tau_0),\\
f_1&=(1,0),&f_2&=(i,0),&f_3&=(1+i,1)/2,&f_4&=(0,\tau_1),
\end{aligned}
\]
where $\rho=e^{2\pi i/3}$ and $\operatorname{Im}\tau_i>0$.
The two rows describe local period coordinates, related by the integral marking
\begin{equation}\label{cc:eq:matrices}
\begin{gathered}
T_0=\begin{pmatrix}0&-1&-1&0\\1&-1&0&0\\0&0&1&0\\0&0&0&1\end{pmatrix},\quad
T_{1,f}=\begin{pmatrix}0&1&0&0\\-1&0&-1&0\\0&0&1&0\\0&0&0&1\end{pmatrix},\\
C=\begin{pmatrix}1&1&0&1\\0&1&0&0\\0&-1&1&0\\0&0&0&1\end{pmatrix},\qquad
T_1=CT_{1,f}C^{-1}.
\end{gathered}
\end{equation}
Columns of $C$ express the $f$ basis in the $e$ basis.

Put $\delta=-2e_1-e_2+3e_3$, $\delta_1=-f_1-f_2+2f_3$, and
$\psi=e_4^*$. Then $C\delta_1=\delta$ and $\psi C=f_4^*$.
All matrices act on column vectors.

The global geometric question concerns a compact complex threefold
$Y(u,a_0,a_1)$ for every sufficiently small nonzero complex parameter $u$,
with primitive torsion increments
\[
a_0\in\tfrac13\Lambda^{T_0},\qquad
a_1\in\tfrac14\Lambda^{T_1},\qquad
|12\psi(a_0+a_1)|=1.
\]
The desired conclusion is a diffeomorphism of its underlying smooth
manifold with $S^6$. The construction below treats the specified pair
\begin{equation}\label{cc:eq:increments}
a_0=e_4/3,\qquad a_1=-Cf_4/4,
\qquad 12\psi(a_0+a_1)=1.
\end{equation}
Here $Cf_4$ means the fourth column of $C$.
The conclusion for all other increments requires their own affine
attachment maps. No quantifier over those increments follows from this pair.

For comparison with complex geometry, assume the following local data at
both finite punctures. The cyclic elliptic reductions have smooth
proper two-torus quotients of the reduced divisor torsors
$\mathcal O(P'_0-O')^\times$. The specified lifted translation extends,
commutes with the deck action, and has the displayed periods.
The original section of $\mathcal O(P_0-O)|_O$ is retained.
Its conormal comparisons are $e=s^{-2}e'$ at $0$ and $e=s^{-1}e'$ at $1$.

These assumptions include the smooth lift and proper quotient.
They do not follow from \eqref{cc:eq:matrices}.
The topological quotients below exist without these geometric assumptions.
Section~\ref{sec:holomorphic-realization} realizes the assumptions
for the selected increments.

\subsection{The common topological diagram}

Set $u_0=e_3$, $u_1=-e_1+e_3$, $u_2=-e_1-e_2+e_3$.
Their oriented three-torus is $X_0$, and $P_0$ cyclically permutes them.
Write $U_0=u_0+u_1+u_2=\delta$.
In the local $f$ coordinates put $X_1=\mathbb R^3/\mathbb Z^3$,
$U_1=(-1,-1,2)$, and
\[
P_1=\begin{pmatrix}0&-1&-1\\1&0&0\\0&0&1\end{pmatrix}.
\]
Thus $P_0^3=P_1^4=I$ and $P_iU_i=U_i$.

For $G_i$ use the convention
$[y,\tau+1]=[P_i y,\tau]$.
The two fillings are
\begin{equation}\label{cc:eq:fillings}
\begin{aligned}
N_0&=(X_0\times S^1_t\times D_s)/
 (x,t,s)\sim(P_0^{-1}x+U_0/3,t+1/3,\rho s),\\
N_1&=(X_1\times S^1_t\times D_s)/
 (x,t,s)\sim(P_1x-U_1/4,t-1/4,is).
\end{aligned}
\end{equation}
The actions are free because their nonidentity powers translate $t$ by
nonintegral multiples of $1/3$ or $1/4$.

Their disk trivializations are
\begin{equation}\label{cc:eq:Phi}
\begin{aligned}
\Phi_0[x,t,s]&=([x+2tU_0,3t],e^{-2\pi it}s),\qquad q=e^{2\pi i\tau}z^3,\\
\Phi_1[x,t,s]&=([x-tU_1,4t],e^{2\pi it}s),\qquad q=e^{-2\pi i\tau}z^4.
\end{aligned}
\end{equation}
Their smooth inverses follow by lifting $\tau$ to $\mathbb R$.
The changes of lift are precisely the quotient relations.
Both shears have determinant one and the factors $3,4$ are positive.
Thus these are orientation-preserving diffeomorphisms to $G_i\times D$.

For $T\in\{T_0,T_1,T_\infty\}$ put
$B_T=(F\times[0,1])/((a,1)\sim(Ta,0))$.
Let $\gamma(v)=[0,v]$ be its positive meridian, based at $[0,0]$.
For a based fiber loop $\lambda(t)$, the square $[\lambda(t),v]$
identifies the two routes from $(0,0)$ to $(1,1)$.
In the declared path order this gives
\begin{equation}\label{cc:eq:square-conjugation}
\gamma\lambda=(T\lambda)\gamma,
\qquad \gamma\lambda\gamma^{-1}=T\lambda.
\end{equation}
The top edge represents $T\lambda$ because $[a,1]=[Ta,0]$.
Thus the endpoint identification and the conjugation formula agree.

The same square in $G_i$ gives $h_i z h_i^{-1}=P_i z$,
where $h_i(v)=[0,v]$ and $z\in\pi_1(X_i,0)$.
The marked boundary maps, in their respective local coordinates, are
\begin{equation}\label{cc:eq:boundary}
\begin{aligned}
j_0[x,t,v]&=[x-2vU_0/3,t+v/3,e^{2\pi iv/3}],\\
j_1[x,t,v]&=[x-vU_1/4,t-v/4,e^{2\pi iv/4}].
\end{aligned}
\end{equation}
The order-four source uses the actual torus isomorphism
$C^{-1}:F\longrightarrow F_1$. In \eqref{cc:eq:Phi} the boundary maps become
\begin{equation}\label{cc:eq:k}
\begin{aligned}
\Phi_0j_0[x,t,v]&=([x+2tU_0,3t+v],e^{-2\pi it}),\\
\Phi_1j_1[x,t,v]&=([x-tU_1,4t-v],e^{2\pi it}).
\end{aligned}
\end{equation}

These formulas include the normal-circle coordinate.
Projection onto $G_i$ gives the actual boundary retractions $k_i$.
Setting $v=0$ gives the actual fiber maps $f_i:F\to G_i$.

The circle matrices on $(t,v)$ are
$\left(\begin{smallmatrix}3&1\\-1&0\end{smallmatrix}\right)$ and
$\left(\begin{smallmatrix}4&-1\\1&0\end{smallmatrix}\right)$.
Both have determinant one. They identify the fiber-first, positive
meridian orientation with the outward-first boundary orientation of $N_i$.

Let $\Sigma$ be the oriented sphere minus the interiors of three disjoint disks.
Construct the flat torus bundle $E\to\Sigma$ with zero section and monodromies
$T_0,T_1,T_\infty=(T_0T_1)^{-1}$.
Concretely, trivialize over a disk obtained by cutting $\Sigma$ along two
arcs, and identify the paired cuts by $T_0$ and $T_1$.
Choose the marking paths along those cuts according to
\eqref{cc:eq:path-order}.

More explicitly, let $\widetilde\Sigma$ consist of endpoint-preserving
homotopy classes of paths $p$ from the base point.
It has the right deck action $p\cdot\gamma=p\gamma$.
Let $\mu:\pi_1(\Sigma)\to\mathrm{SL}(\Lambda)$ send
$\gamma_i$ to $T_i$ for $i=0,1$.
Define the bundle as the quotient by the right action
\[
(a,p)\cdot\gamma=(\mu(\gamma)^{-1}a,p\gamma).
\]
Indeed two successive actions by $\alpha,\beta$ give
$\mu(\beta)^{-1}\mu(\alpha)^{-1}=\mu(\alpha\beta)^{-1}$.

Let $p_*$ denote the constant path at the base point.
A lifted loop ends at $[a,p_*\gamma]=[\mu(\gamma)a,p_*]$.
Hence its mapping torus has exactly the stated endpoint identification.
Transport satisfies $\mu(\alpha\beta)=\mu(\alpha)\mu(\beta)$.
Equation~\eqref{cc:eq:path-order} therefore gives
$T_0T_1T_\infty=I$.

Its boundary collars are the specified $B_{T_i}$.
Since the boundary orientation of $\Sigma$ is clockwise around each
removed disk, the maps $j_i$ reverse the induced boundary orientations
for gluing. Define the compact oriented manifold with boundary
\begin{equation}\label{cc:eq:W}
W=E\mathop{\cup}_{j_0}N_0\mathop{\cup}_{j_1}N_1,
\qquad \partial W=B_{T_\infty}
\end{equation}
with the inherited boundary orientation.

Comparison with the original complex family requires identifying its
geometric torus bundle with this flat bundle.
The identification must preserve all three chosen boundary sections.

\subsection{A finite integral chain model}

We specify chains by oriented simplices of the actual spaces.
This construction permits subdivisions between different fiber sections.
It does not require a triangulation invariant under the whole monodromy group.

Give each torus a rational polyhedral structure in its displayed real
period coordinates. For $X_i$, refine a sufficiently fine cubic tiling by
its finitely many images under $P_i$, and subdivide barycentrically.
This produces an actual finite triangulation on which $P_i$ acts.
Use its mapping cylinder, with interval-first prisms, for $G_i$.
Triangulate the disks as polygons with a central vertex.

The maps \eqref{cc:eq:k} are piecewise linear in the circle coordinates
$(t,v)$ after cutting at integer values of $3t+v$ and $4t-v$.
Cut further at inverse images of the finitely many target faces.
Subdivide these rational polytopes and their faces compatibly.
This supplies triangulations of the boundary and marked fiber subspaces
of $G_i\times D$, and triangulates their inclusions.
Here circle arguments, rather than Cartesian disk coordinates, describe
the boundary. Extension to the disk is by coning its polygon.

Transport these triangulations to $\partial E$ using \eqref{cc:eq:k} and $C$.
Extend them across collars and then across $E$.
Cut the base to the polygon used above.
Subdivide both sides of each cut by the transported torus triangulations,
then interpolate between these subdivisions in a collar.
For two rational torus subdivisions, their intersection is a finite
polyhedral subdivision. Triangulate the product collar using that
common subdivision at its middle level.

The remaining product over the cut polygon has a product triangulation.
At the corners the transition product is the identity, so the
identified faces agree. This constructs a finite triangulation of the
bundle with the prescribed boundary triangulations.
Further collar subdivisions make the gluing maps simplicial.

Let $A,L_i,B_i$ denote the resulting simplicial integral chain complexes
of $E,N_i,B_{T_i}$, respectively, for $i=0,1$.
Let $\iota_i:B_i\to A$ and $j_i:B_i\to L_i$ be their inclusion maps.

This prescription defines every integer entry: an oriented simplex has
boundary
\[\sum_{r=0}^n(-1)^r[v_0,\ldots,\widehat v_r,\ldots,v_n].\]
Each inclusion sends it to the same oriented simplex in the larger space.
If the chosen orientations disagree, its coefficient is $-1$.
The changes of periods and the angular maps have already been used to
identify the simplex, before assigning its coefficient.
Thus the construction gives a finite chain diagram with genuine maps,
including all dimensions through six. Enumerating its incidence matrices
is optional for its definition, but necessary for a numerical Smith computation.

\begin{theorem}\label{cc:thm:cone}
Put $B=B_0\oplus B_1$, $A'=A\oplus L_0\oplus L_1$, and
\[
J(b_0,b_1)=(\iota_0b_0+\iota_1b_1,-j_0b_0,-j_1b_1),
\]
where $\iota_i$ denotes the inclusion into $A$.
The finite chain complex
\begin{equation}\label{cc:eq:cone}
\mathcal W_n=A'_n\oplus B_{n-1},\qquad
d(a,b)=(d a+Jb,-d b)
\end{equation}
computes $H_*(W;\mathbb Z)$.
The comparison sends $(a,b)$ to the sum of the chains $a$ in $W$.
Its boundary component has image zero.
For the inclusion $b_\infty:C_*(B_{T_\infty})\to A$,
the actual cusp boundary map into \eqref{cc:eq:cone} is
\begin{equation}\label{cc:eq:cusp-left}
\beta(b)=(b_\infty b,0,0;0,0).
\end{equation}
All coefficients and comparisons are integral.
\end{theorem}

\begin{proof}
The map $J$ is injective because each boundary inclusion into its filling
is injective on simplices. Its cokernel is exactly $C_*(W;\mathbb Z)$.
The kernel of the stated cone comparison consists of pairs $(Jc,b)$.
Its differential is $(c,b)\mapsto(dc+b,-db)$ in these coordinates.
The degree-one homotopy $(c,b)\mapsto(0,c)$ contracts it.
This proves the comparison and also the usual integral Mayer--Vietoris
sequence.

Formula \eqref{cc:eq:cusp-left} is a chain map because the
cusp boundary is a subcomplex disjoint from the two glued boundaries.
\end{proof}

The fiber inclusions $F\to B_{T_i}\to N_i$ in this diagram induce the
specialization maps by dualizing integer chains.
Their deformation to the core is exactly \eqref{cc:eq:k} at $v=0$.
\subsubsection{Local prism complexes and their maps}

Put $D_i=C_*^{\mathrm{simp}}(X_i;\mathbb Z)$ for the chosen
$P_i$-invariant triangulation, and write $p_i=(P_i)_*$.
Let $D_{i,n}=0$ outside $0\leq n\leq3$.
The symbol $p_i$ denotes the induced chain automorphism in every degree.
It commutes with $\partial$ and satisfies $p_0^3=p_1^4=I$.

Use the product cells of $X_i\times S^1$ with one vertex and one
positive edge on the circle. Its chain complex is
\begin{equation}\label{cc:eq:product-carrier}
Q_{i,n}=D_{i,n}\oplus D_{i,n-1},\qquad
d_Q(x,y)=(\partial x,-\partial y).
\end{equation}
Here $x$ is a fiber chain and $y$ represents the interval-first cell
$[0,1]\times y$.

The mapping-torus cells of $G_i$ give
\begin{equation}\label{cc:eq:torus-carrier}
M_{i,n}=D_{i,n}\oplus D_{i,n-1},\qquad
d_M(x,y)=(\partial x+(p_i-I)y,-\partial y).
\end{equation}
The incidence term follows from
\[
\partial([0,1]\times y)=\{1\}\times y-\{0\}\times y
 -[0,1]\times\partial y.
\]
The two endpoint chains become $p_i y$ and $y$ under the quotient.
This also fixes the sign of the shifted differential.

The fourth periods in these product coordinates are
$\bar w_0=e_4-2\delta$ and $\bar w_1=f_4+\delta_1$.
The shears in \eqref{cc:eq:k} identify the fiber maps with
$[y,t]\mapsto[y,d_i t]$, where $(d_0,d_1)=(3,4)$.
For instance $x=y-2tU_0$ at $0$, and $x=y+tU_1$ at $1$.
These expressions send a positive fourth period to $\bar w_i$.

The source monodromies on $Q_i$ are
\[
\mathcal T_0(x,y)=(p_0x,p_0y),\qquad
\mathcal T_1(x,y)=(p_1^{-1}x,p_1^{-1}y).
\]
Thus $\mathcal T_1$ is the actual order-four monodromy in the local
$f$ marking. The change $C$ carries it to $T_1$ in the common marking.

\begin{lemma}[Fiber maps and meridian prisms]\label{cc:lem:prisms}
Put $\nu_i=I+p_i+\cdots+p_i^{d_i-1}$.
The fiber maps and positive meridian prisms induce
\begin{equation}\label{cc:eq:local-maps}
\begin{aligned}
F_i:Q_{i,n}&\longrightarrow M_{i,n},& F_i(x,y)&=(x,\nu_i y),\\
H_0:Q_{0,n}&\longrightarrow M_{0,n+1},&H_0(x,y)&=(0,x),\\
H_1:Q_{1,n}&\longrightarrow M_{1,n+1},&H_1(x,y)&=(0,-p_1^{-1}x).
\end{aligned}
\end{equation}
They satisfy
\begin{equation}\label{cc:eq:local-homotopy}
d_MF_i=F_id_Q,\qquad
d_MH_i+H_id_Q=F_i(\mathcal T_i-I).
\end{equation}
In the boundary mapping-torus complex
\[
R_{i,n}=Q_{i,n}\oplus Q_{i,n-1},\qquad
d_R(a,b)=(d_Qa+(\mathcal T_i-I)b,-d_Qb),
\]
the boundary retraction is the chain map
\begin{equation}\label{cc:eq:boundary-prism-map}
\mathcal K_i:R_{i,n}\longrightarrow M_{i,n},\qquad
\mathcal K_i(a,b)=F_i a+H_i b.
\end{equation}
Subdivision and the deformation retractions identify this map with
the actual boundary maps in the finite attachment diagram.
\end{lemma}

\begin{proof}
Divide the fourth interval into $d_i$ equal parts.
Their images traverse the successive core intervals with coefficients
$I,p_i,\ldots,p_i^{d_i-1}$. This gives $F_i$.

In the adapted coordinates, the two boundary retractions are
\[
[y,t,v]\longmapsto[y,3t+v],\qquad
[y,t,v]\longmapsto[y,4t-v].
\]
A positive $v$ interval therefore gives one positive core interval at $0$.
At $1$, it runs from core parameter $0$ to $-1$.
The identity $[y,\tau-1]=[P_1^{-1}y,\tau]$ gives coefficient
$-p_1^{-1}$. These are the displayed $H_i$ on the fiber components.

A cell containing both the $t$ and $v$ intervals maps through one core
coordinate, so its top cellular coefficient is zero.
This proves the formula on the second components.
The maps are cellular for the product prism structures.
Subdividing their prisms gives the same induced maps on simplicial chains,
up to the chain homotopies of the subdivision comparisons.

Algebraically, $(p_i-I)\nu_i=\nu_i(p_i-I)=0$ proves that $F_i$ is a
chain map. Direct substitution gives
\[
\begin{aligned}
(d_MH_0+H_0d_Q)(x,y)&=((p_0-I)x,0),\\
(d_MH_1+H_1d_Q)(x,y)&=((p_1^{-1}-I)x,0).
\end{aligned}
\]
The second components of $F_i(\mathcal T_i-I)$ vanish because
$\nu_i(p_i^{\pm1}-I)=0$.
Finally,
\[
d_M\mathcal K_i(a,b)
=F_id_Qa+F_i(\mathcal T_i-I)b-H_id_Qb
=\mathcal K_id_R(a,b).
\]
Thus the same positive interval convention determines the endpoint
incidence, the meridian conjugation, and the boundary chain map.
\end{proof}

The integral filling homology is free of ranks $(1,2,2,2,1)$ in each case.
The cokernels of $F_{0*}$ in degrees zero through four are
$0,\mathbb Z/3,0,0,\mathbb Z/3$.
For $F_{1*}$ they are
$0,\mathbb Z/4,\mathbb Z/2,\mathbb Z/2,\mathbb Z/4$.

To verify the integral assertions, use the exact sequence of
\eqref{cc:eq:torus-carrier}:
\[
0\longrightarrow\operatorname{coker}(P_i-I\mid H_n(X_i))
\longrightarrow H_n(G_i)
\longrightarrow\ker(P_i-I\mid H_{n-1}(X_i))\longrightarrow0.
\]
Here and below homology has integral coefficients.
The fiber component of $F_{i*}$ surjects onto the left term.
Its fourth-interval component induces the norm on the right term.
Consequently
\[
\operatorname{coker}F_{i*,n}
\cong
\frac{\ker(P_i-I\mid H_{n-1}(X_i))}
{(I+P_i+\cdots+P_i^{d_i-1})H_{n-1}(X_i)}.
\]

For $P_0$, the degree-one and degree-two actions are cyclic permutations.
Their norm images generate the invariant lattices.
For $P_1$, the degree-one and degree-two matrices are
\[
\begin{pmatrix}0&-1&-1\\1&0&0\\0&0&1\end{pmatrix},\qquad
\begin{pmatrix}1&1&0\\0&0&-1\\0&1&0\end{pmatrix}.
\]
Their differences from the identity have Smith factors $(1,1)$.
Their norm images are $2\mathbb Z\delta_1$ and
$2\mathbb Z(f_1\wedge f_2)$.

In degrees zero and three the norm is
multiplication by the covering degree. These facts prove the assertions,
including the possible finite cokernels. They do not compute the homology
of the global cone.

\begin{proposition}\label{cc:prop:no-common}
There is no finite triangulation of $F$ on which both actual maps
$T_0,T_1$ are simplicial automorphisms.
\end{proposition}

\begin{proof}
Multiplication in the common marking gives
\begin{equation}\label{cc:eq:cusp-monodromy}
T_\infty=\begin{pmatrix}2&1&1&0\\0&1&0&-1\\-1&-1&0&1\\0&0&0&1\end{pmatrix}
=I+N,\qquad N^2=0,\quad N\ne0.
\end{equation}
Thus $T_\infty^r=I+rN$ has infinite order over $\mathbb Z$.
Simplicial automorphisms of a finite triangulation form a finite group,
since they permute its vertices faithfully. Its induced action on
$H_1(F;\mathbb Z)$ is therefore finite, a contradiction.
\end{proof}

This obstruction concerns invariant triangulations. It does not obstruct
the finite attachment diagram or the compact manifold $W$.

\subsection{Meridians in one integral marking}

Let $h_i$ be the positive core interval of $G_i$, as in
\eqref{cc:eq:square-conjugation}.
The formulas \eqref{cc:eq:k} give $\gamma_0\mapsto h_0$ and
$\gamma_1\mapsto h_1^{-1}$.
Consequently the based relations are
\begin{equation}\label{cc:eq:meridians}
\gamma_0^3=f_{0\#}(w_0),\qquad
\gamma_1^4=f_{1\#}(w_1),\qquad
w_0=e_4-2\delta,\quad w_1=-Cf_4-\delta.
\end{equation}
In $B_0$ coordinates, $w_0=(4,2,-6,1)^t$ and
$w_1=(1,1,-3,-1)^t$.
The second equation is exactly
$\gamma_1^4=f_{1\#}(-f_4-\delta_1)$ in its local marking.

The fiber maps are coverings of degrees three and four onto the cores.
Their induced fundamental-group maps are injective.
Thus the nonzero conormal terms cannot be erased in either local group.

Indeed the product coordinates above give
$f_{0\#}(\bar w_0)=h_0^3$ and $f_{1\#}(\bar w_1)=h_1^4$.
The marked meridians are $h_0$ and $h_1^{-1}$, so
$w_0=\bar w_0$ and $w_1=-C\bar w_1$.
For $z\in\pi_1(X_i)$, the same square gives conjugation by $P_i$
for $h_i$, and by $P_i^{-1}$ for $h_i^{-1}$.
This recovers $T_0$ and $T_{1,f}$, including the order-four inverse.

The fourth generators commute with their meridians because $P_i^{d_i}=I$.
Hence each boundary group has presentation
\[
\langle\Lambda,\gamma_i\mid[\Lambda,\Lambda]=1,
\ \gamma_i\lambda\gamma_i^{-1}=T_i\lambda\rangle,
\]
and its map to the core adds exactly the power relation in
\eqref{cc:eq:meridians}. The covering subgroup is $\Lambda$.

The restriction of a degree-one integral class $\alpha$ on the fiber
extends to the local core exactly when
\begin{equation}\label{cc:eq:congruences}
\alpha T_i=\alpha,\qquad d_i\mid\alpha(w_i),\qquad(d_0,d_1)=(3,4).
\end{equation}
Indeed a homomorphism from the core group to $\mathbb Z$ must assign
$\alpha(w_i)/d_i$ to the marked meridian.
Conversely this assignment, together with $\alpha$ on the fiber subgroup,
respects its conjugation and power relations and defines a homomorphism.

The only common invariant covectors are $\mathbb Z\psi$.
For $r\psi$, the two divisibility conditions are $3\mid r$ and $4\mid r$.
Their intersection is $12\mathbb Z\psi$.

In particular the local core extensions of $12\psi$ evaluate to $4$ and
$-3$ on the two positive meridians. These evaluations include the
conormal terms, since $\psi(\delta)=0$.
They do not yet specify an extension over the cusp.

\begin{proposition}\label{cc:prop:pi1}
With additive notation for the commuting fiber subgroup,
\[
\pi_1(W)=\langle\Lambda,\gamma_0,\gamma_1\mid
 [\Lambda,\Lambda]=1,\quad
 \gamma_i\lambda\gamma_i^{-1}=T_i\lambda,\quad
 \gamma_0^3=w_0,\ \gamma_1^4=w_1\rangle.
\]
Moreover $H_1(W;\mathbb Z)\cong\mathbb Z$.
One primitive generator $g$ satisfies
$[e_4]=12g$, $[\gamma_0]=4g$, and $[\gamma_1]=-3g$.
In particular $[\gamma_\infty]=-g$.
\end{proposition}

\begin{proof}
The flat bundle retracts onto its restriction to a bouquet of two circles.
Its zero section splits its fundamental-group extension.
Formula~\eqref{cc:eq:square-conjugation} gives its conjugation action,
with multiplication as in \eqref{cc:eq:path-order}.
Each boundary group maps onto its core group. Its kernel is normally
generated by the corresponding relation \eqref{cc:eq:meridians}, as
the fiber is the index-$d_i$ covering subgroup. Van Kampen gives the presentation.

For abelianization, the images of $T_0-I$ contain $-e_1$ and $e_2-e_1$.
The image of $T_1-I$ contains $-2e_1-e_2+e_3$.
They therefore generate $\ker\psi$ integrally.

The remaining relations are $3x_0=c$ and $4x_1=-c$.
Their relation matrix on $(c,x_0,x_1)$ is
$\left(\begin{smallmatrix}-1&3&0\\1&0&4\end{smallmatrix}\right)$.
Its maximal minors have absolute values $3,4,12$, with gcd one.
Thus there is no torsion and the primitive solution is $(c,x_0,x_1)=(12,4,-3)g$.
The infinity relation gives
$[\gamma_\infty]=-[\gamma_0]-[\gamma_1]=-g$.
\end{proof}

\subsection{An explicit cusp collapse and its obstruction}

The primitive lattice of cusp vanishing cycles is
\[
K=N\Lambda=\langle e_1-e_3,e_3-e_2\rangle.
\]
Define $A_\infty:\Lambda\to\mathbb Z^2$ by
\begin{equation}\label{cc:eq:A}
A_\infty=\begin{pmatrix}1&1&1&0\\0&0&0&1\end{pmatrix}.
\end{equation}
It is surjective, has kernel $K$, and satisfies $A_\infty T_\infty=A_\infty$.
It therefore defines the actual map
\begin{equation}\label{cc:eq:p}
p_\infty:B_{T_\infty}\longrightarrow\mathbb R^2/\mathbb Z^2,
\qquad [a,v]\longmapsto A_\infty a.
\end{equation}
It kills the marked meridian and the two vanishing cycles.

Let $Z$ be its mapping cylinder and attach its boundary end to $W$.
The resulting finite polyhedron is denoted $W_p$.
All maps here are explicit piecewise linear maps after finite subdivisions.

\begin{theorem}\label{cc:thm:cusp-obstruction}
The space $W_p$ is simply connected and has Euler characteristic zero.
It is not an integral homology six-sphere.
Consequently the vanishing-cycle quotient and a null-homotopic cusp
meridian do not determine the cusp chains required by that conclusion.
\end{theorem}

\begin{proof}
The map $p_\infty$ is surjective on fundamental groups.
Its kernel is normally generated by $K$ and $\gamma_\infty$.
The normal closure of $K$ in $\pi_1(W)$ contains
\[
e_1-e_3,\qquad e_3-e_2,\qquad T_0(e_1-e_3)=e_1+e_2-e_3.
\]
Their first three coordinates have determinant $-1$. They generate
$\ker\psi$.

Equation~\eqref{cc:eq:path-order} gives
$\gamma_\infty=\gamma_1^{-1}\gamma_0^{-1}$.
Killing this meridian therefore imposes $\gamma_0\gamma_1=1$.
It has the presentation
\[
\langle c,x\mid [c,x]=1,\ x^3=c,\ x^{-4}=c^{-1}\rangle=1.
\]
This proves simple connectivity by van Kampen.

The flat torus bundle has Euler characteristic
$\chi(E)=\chi(\Sigma)\chi(F)=0$.
Each $N_i$ retracts onto a mapping torus, and has Euler characteristic zero.
Each boundary mapping torus also has Euler characteristic zero.
Additivity therefore gives $\chi(W)=0$.
The mapping cylinder $Z$ retracts onto a two-torus, so $\chi(Z)=0$.
Hence $\chi(W_p)=0+0-0=0$, whereas an integral homology six-sphere
has Euler characteristic two.
\end{proof}

The mapping cylinder is a topological comparison, and is not asserted
to be a smooth cusp neighborhood. Its failure is already detected
integrally by Euler characteristic. Toric strata of an actual cusp can
supply additional chains even when their fundamental-group effect
agrees with \eqref{cc:eq:p}.

\subsection{The cusp map required by the common cone}

Let $N_\infty\to D$ be a compact oriented six-manifold with a specified
boundary diffeomorphism $j_\infty:B_{T_\infty}\to\partial N_\infty$.
It must extend the prescribed torus family and the chosen cusp section.
For complex-geometric application, require a holomorphic extension with
smooth total space and the prescribed punctured identification.
Its boundary orientation must be opposite to that of $W$.

After triangulating the pair, put $L_\infty=C_*(N_\infty;\mathbb Z)$ and
$\kappa_\infty=(j_\infty)_*$, including the boundary into $N_\infty$.
This denotes an induced chain map only when the map $j_\infty$ is supplied.
It is not determined by \eqref{cc:eq:A}.

\begin{theorem}[Conditional integral gluing]\label{cc:thm:conditional}
For these geometric data, the union $Y=W\cup_{j_\infty}N_\infty$ has integral
chains represented by
\begin{equation}\label{cc:eq:globalcone}
\begin{split}
\mathcal Y_n&=\mathcal W_n\oplus L_{\infty,n}\oplus C_{n-1}(B_{T_\infty}),\\
d(w,l,b)&=(dw+\beta b,dl-\kappa_\infty b,-db).
\end{split}
\end{equation}
It has the integral homology of a six-sphere precisely when this cone
has homology $\mathbb Z$ in degrees zero and six, and zero elsewhere.
For an explicit criterion, choose an integral basis of $\ker d_n$ in
each degree, and express $\operatorname{im}d_{n+1}$ in that basis.
Every nonzero Smith factor of this presentation must be one.
Its cokernel must have rank one in degrees zero and six, and rank zero
in all other degrees.
The claim includes the comparison of \eqref{cc:eq:globalcone} with the
specified attachment, not just a list of its ranks over a field.
\end{theorem}

\begin{proof}
Let $q:\mathcal W\to C_*(W;\mathbb Z)$ be the comparison in
Theorem~\ref{cc:thm:cone}.
It satisfies $q\beta=(B_{T_\infty}\hookrightarrow W)_*$ exactly.
Consequently $(w,l,b)\mapsto(qw,l,b)$ maps \eqref{cc:eq:globalcone}
onto the ordinary attachment cone for $W$ and $N_\infty$.
Its kernel is the contractible complex $\ker q$ described there.

The ordinary attachment cone maps onto $C_*(Y;\mathbb Z)$.
Its kernel is again the identity cone, with contraction
$(c,b)\mapsto(0,c)$.
The composite is therefore an integral quasi-isomorphism.

Each cycle group is a subgroup of a finite free abelian group, hence
has a finite integral basis.
The stated Smith calculation computes the quotient of cycles by boundaries,
which is exactly the required integral homology.
\end{proof}

For a deformation retraction $r_\infty:N_\infty\to Y_\infty$ onto a cusp
central fiber, the deciding specialization is the concrete composite
\begin{equation}\label{cc:eq:deciding}
F\longrightarrow B_{T_\infty}
 \mathrel{\mathop{\longrightarrow}^{j_\infty}}N_\infty
 \mathrel{\mathop{\longrightarrow}^{r_\infty}}Y_\infty.
\end{equation}
One must give its maps on the toric cells and their intersections,
including the images of the marked fiber and meridian prisms.
The chain homotopy between inclusion and this retraction supplies the
comparison with \eqref{cc:eq:globalcone}.
In particular the degree-two and degree-three specialization images
cannot be inferred from $K$ or from $\pi_1(Y)$.
For any homology-sphere attachment here, Euler additivity requires
$\chi(N_\infty)=2$.

The finite fillings provide positive normal disks.
In a boundary cone with differential $(a,b)\mapsto(da+kb,-db)$, let
$Lq$ insert the positive normal-circle chain first.
Then $dL=-Ld$ and $kL=0$ on normalized cubical chains.
The positive Thom map is $q\mapsto(0,-Lq)$.
Indeed the radial homotopy sends its relative image to
$(r,v)\mapsto re^{2\pi iv}$ times $q$, which is the positive disk class.
This fixes the cone sign independently of the direction of the core loop.

The cusp comparison must preserve the induced orientation at its own
boundary. An asserted scalar transgression does not supply that comparison.

The toric construction in Section~\ref{sec:toric-cusp} supplies the
central surface, and Section~\ref{sec:cusp-attachment} constructs
the map \eqref{cc:eq:deciding}. The collapse theorem explains why
the primitive vanishing-cycle quotient cannot replace that surface.
Comparison with a complex family also requires identifying the regular
bundles relative to their sections.
This identification must extend over the finite fillings. Section~\ref{sec:geometric-comparison}
constructs these maps for the selected lifts and adapted cusp. Their
geometric construction is distinct from the integral calculation for
the common cone.
